\protect \section *{\nameref *{sec:plasma-kinetic}}
\begin{Solution}{7.{1}}
		$f\to0$ при $|\vect v|\to\infty$ швидше за~$1/v^2$; крім того, для сили Лоренца $\nabla_{\vect v}\cdot\vect F=0$.
	
\end{Solution}
\begin{Solution}{7.{2}}
		для ізотропного $f$ тензор $\mathsf P_{ij}=p\,\delta_{ij}$, тож $p=\operatorname{Tr}\mathsf P/3$; для максвеллівського розподілу $\int c^2f\,d^3v=3nT/m$.
	
\end{Solution}
\begin{Solution}{7.{3}}
		(а)~$\lambda_{De}\approx2{,}4\times10^{-3}$~см, $k\lambda_{De}\approx0{,}15$, $\omega\approx1{,}03\,\omega_{pe}$,
$v_\varphi\approx9\times10^8$~см/с, $v_\varphi/v_{Te}\approx7$, $\gamma/\omega_{pe}\approx5\times10^{-9}$~--- загасання дуже слабке;
(б)~$\gamma/\omega\approx1/(2\pi\cdot10)\approx0{,}016$ дає $k\lambda_{De}\approx0{,}3$, тобто $\lambda\approx2\pi\lambda_{De}/0{,}3\approx0{,}5$~мм.
	
\end{Solution}
\begin{Solution}{7.{4}}
		$\sqrt{2T/m}$, $\sqrt{8T/\pi m}$, $\sqrt{3T/m}$
(у~відношенні $1:1{,}13:1{,}22$).
	
\end{Solution}
\begin{Solution}{7.{5}}
		$\delta f=-i\,(q/m)\,E\,F'/(\omega-kv+i\nu)$, тобто $\omega\to\omega+i\nu$.
	
\end{Solution}
\begin{Solution}{7.{7}}
		$1=\dfrac{\omega_p^2}2\Bigl[\dfrac1{(\omega-ku)^2}+\dfrac1{(\omega+ku)^2}\Bigr]$.
	
\end{Solution}
\begin{Solution}{7.{8}}
		$\omega_b\approx0{,}01\,\omega_{pe}$, $\gamma\approx0{,}02\,\omega_{pe}$: хвиля загасає швидше, ніж
починається захоплення.
	
\end{Solution}
\begin{Solution}{7.{9}}
		$F_b'\sim n_b/\Delta v^2$ максимальна на~висхідному схилі пучка, при $v\approx v_b-\Delta v$.
	
\end{Solution}
\begin{Solution}{7.{10}}
		$T_e=m_ec^2\bigl[\sigma_\lambda/(2\lambda_i\sin45^\circ)\bigr]^2\approx90$~еВ,
$\lambda_{De}\approx10^{-3}$~см, $\alpha\approx0{,}01\ll1$.
	
\end{Solution}
\begin{Solution}{7.{11}}
		$\sin(\theta/2)=\lambda_i/4\pi\lambda_{De}\approx0{,}11$, $\theta\approx13^\circ$; колективне при $\theta\lesssim13^\circ$.
Для 1064~нм потрібно $\sin(\theta/2)\approx0{,}011$, тобто $\theta\approx1^\circ$~--- відокремити сигнал від~пучка практично неможливо.
	
\end{Solution}
